\chapter{スキーム論}
スキーム論の試問対策を作っていく. Liの個人的なノートでもある. 
\section{いろいろな基礎}



\subsection{スペクトラム}

\begin{egprob} 
$\spec{A}$が1点のとき$A$は体ですか?
\end{egprob}
\begin{egans}
No. $A=\mb{C}[x]/(x^2)$.
\end{egans}


\subsection{スキーム}
\begin{egprob}
アファインでないスキームの例は? 
\end{egprob}
\begin{egans}
(\tf{例1}) $X = \spec{k[x,y]} \setminus{\set{(x,y)}}$. まずこれは整域のスペクトラム. $\mac{O}_{X}(X)$は$\spec{k[x,y]}$の$D(x)\cup D(y)$なので$k[x,y]_{x} \cap k[x,y]_{y} = k[x,y,\dfrac{1}{x}] \cap k[x,y,\dfrac{1}{y}] = k[x,y]$. よって$X$と$\spec{k[x,y]}$は同じ大域断面を持つ. もし$X$がアファインなら$\mr{id}_{k[x,y]}$に対応する$\spec{k[x,y]}\to X$という射がある. これが同型射であるべきなので$X\cong \spec{k[x,y]}$で, とくに開埋入$X\to \spec{k[x,y]}$が全単射だがそれはありえない.  \\
(\tf{例2}) $n\geq 1$に対し射影空間$\mb{P}^{n}_{k}$. 大域断面が$k$なので$\spec{k}$になるべきだが明らかに点の個数が違う. 
\end{egans}

\begin{egprob}

\end{egprob}




























\newpage 
\section{レポート対策！}


\begin{egprob}
Noether正規化補題の証明のアウトラインを教えてください.
\end{egprob}

\begin{egprob}
強零点定理については知っていますか? 
\end{egprob}











